\section{The basics, from zero}\label{sec:basics}

\subsection{Numbers inside a computer are bits}

A computer stores every number as a row of \emph{bits}. A bit is a tiny switch that is either
off (0) or on (1). Just as our everyday numbers use the place values 1, 10, 100, 1000 (each place
worth ten times the one to its right), binary numbers use the place values 1, 2, 4, 8, 16, \dots
(each place worth \emph{two} times the one to its right).

\begin{figure}[H]
\centering
\begin{tikzpicture}[box/.style={draw,minimum width=1.05cm,minimum height=0.8cm,font=\large\ttfamily}]
  \foreach \v [count=\i] in {128,64,32,16,8,4,2,1} {
    \node[font=\small,gray] at (\i*1.1,1.0) {\v};
  }
  \foreach \b [count=\i] in {0,0,0,0,1,1,0,1} {
    \ifnum\b=1
      \node[box,fill=roba!20] at (\i*1.1,0) {\b};
    \else
      \node[box] at (\i*1.1,0) {\b};
    \fi
  }
  \node[anchor=west] at (9.4,0) {$= 8+4+1 = 13$};
  \node[anchor=east,font=\small,gray] at (0.4,1.0) {place value};
  \node[anchor=east,font=\small] at (0.4,0) {bits of 13};
\end{tikzpicture}
\caption{The number 13 as eight bits: \code{00001101}. The bits that are 1 sit under 8, 4 and 1,
and $8+4+1=13$.}
\end{figure}

A few more examples:
\[
  11 = \texttt{1011} = 8+2+1,\qquad 100 = \texttt{1100100} = 64+32+4,\qquad
  200 = \texttt{11001000} = 128+64+8.
\]

\paragraph{How many bits?} With 8 bits you can store $0$ to $255$; with one bit used for the sign,
$-127$ to $127$ (that is the \textbf{int8} type used by our networks). With 16 bits you can go up
to about $\pm 32{,}767$ (\textbf{int16}); with 32 bits up to about $\pm 2$ billion (\textbf{int32}).

\subsection{Shifting: the cheapest way to double}

If you write a 0 at the end of an everyday number, it becomes ten times bigger:
$13 \to 130$. In binary, writing a 0 at the end makes the number \emph{two} times bigger. Moving
all bits one place to the left is called a \textbf{left shift}, written \code{<<}.
\[
  13 = \texttt{1101}\quad\xrightarrow{\;<<1\;}\quad \texttt{11010} = 26,\qquad
  13 \ll 3 = \texttt{1101000} = 104 = 13\times 8.
\]
Shifting left by $k$ places multiplies by $2^k$. Shifting \emph{right} (\code{>>}) divides by
$2^k$ and throws away the remainder: $13 \gg 1 = \texttt{110} = 6$.

\begin{keyidea}
For a CPU, a shift costs \textbf{one clock cycle}, exactly like an addition. So multiplying by
1, 2, 4, 8, 16, \dots{} (a \emph{power of two}) is basically free. Multiplying by any other
number, like 11, is the hard part. All the tricks in this project are about turning hard
multiplications into shifts and additions.
\end{keyidea}

\subsection{Multiplying with only additions and shifts}\label{sec:shiftadd}

School multiplication works digit by digit. Binary multiplication works the same way, but every
digit is 0 or 1, so each step is either ``add a shifted copy'' or ``add nothing''.

\begin{example}[: $13\times 11$ by shift-and-add]
$11 = \texttt{1011}$. Look at the bits of 11 from the right. Each time, double 13; when the bit is 1,
add the current copy of 13 to the result.

\begin{center}
\begin{tabular}{cccccc}
\toprule
Step & bits of $b$ left & lowest bit & copy of $a$ & add? & result so far\\
\midrule
1 & \texttt{1011} & 1 & 13 & yes & 13\\
2 & \texttt{101}  & 1 & 26 & yes & 39\\
3 & \texttt{10}   & 0 & 52 & no  & 39\\
4 & \texttt{1}    & 1 & 104 & yes & \textbf{143}\\
\bottomrule
\end{tabular}
\end{center}
Indeed $13\times 11 = 143$. It took 4 rounds because 11 has 4 bits. An 8-bit number needs up to 8
rounds, a 16-bit number up to 16, a 31-bit number up to 31.
\end{example}

In C, the loop is only a few lines (this is Listing~1 of the paper, and the routine
\code{roba\_shift\_add} in our \file{csrc/roba\_core.h}):

\begin{lstlisting}[language=C]
uint32_t roba_shift_add(uint32_t a, uint32_t b)
{
    uint32_t r = 0;
    while (b) {           /* until no bits of b are left        */
        if (b & 1u)       /* lowest bit of b is 1?              */
            r += a;       /*   then add the current copy of a   */
        a <<= 1;          /* double a                           */
        b >>= 1;          /* drop the lowest bit of b           */
    }
    return r;
}
\end{lstlisting}

Each round runs about 6 instructions (test, branch, maybe add, two shifts, loop back), so an
exact 8-bit multiply costs tens of clock cycles, and it gets slower as numbers get bigger.
Our measurements later (Section~\ref{sec:multbench}) show 84 cycles for an 8-bit exact multiply
and 280 cycles for a 31-bit one.

\subsection{What a CPU is, and what a ``cycle'' is}

A \textbf{CPU} (processor) is the part of a chip that runs a program. A program is a long list of
very simple \textbf{instructions}, such as ``add these two numbers'', ``shift this number left by 3'',
``load a byte from memory'', ``if this is zero, jump over there''. The CPU keeps a few dozen
numbers close at hand in \textbf{registers} (32 of them on RISC-V, named \code{a0}, \code{a1},
\code{t0}, \dots).

The CPU is driven by a \textbf{clock}: a signal that ticks millions of times per second. One tick is
one \textbf{clock cycle}. A simple instruction like \code{add} takes one cycle; loading from memory
or a taken jump may take two or three.

\begin{example}[: from cycles to seconds]
If the chip's clock runs at 50~MHz (50 million ticks per second, a typical speed for a small
RISC-V core on an FPGA board), then one cycle lasts $1/50{,}000{,}000$~s $=20$~nanoseconds.
Our exact-multiply MLP needs 13.5 million cycles per image, i.e.\ $13.5\times 10^6 \times 20\,
\text{ns} = 0.27$~s per picture. With RoBA-Dn1 (7.05 million cycles) it is 0.14~s.
\end{example}

Because the clock speed depends on the chip, we always report \textbf{cycles}, not seconds.
Fewer cycles means faster \emph{and} (at the same power) less energy.

\subsection{RISC-V, Ibex, and why some chips have no multiplier}

\textbf{RISC-V} is a free, open \emph{instruction set}: a published list of which instructions a CPU
must understand. The base list, \textbf{RV32I}, has additions, shifts, logic, loads, stores and
jumps, but \textbf{no multiplication}. Multiplication is an optional add-on called the
\textbf{M extension}. Another add-on, \textbf{C}, adds short 16-bit versions of common
instructions to make programs smaller. Our CPU is \textbf{RV32IC}: base + compressed, no M.

\textbf{Ibex} is a small, real, open-source RISC-V CPU made by the company lowRISC. It is used in
real chips (for example in security chips) and can be built with or without a multiplier.

Why would anyone leave out the multiplier? A hardware multiplier is a big block of logic. The
paper synthesised Ibex for a 65~nm chip process and found that the multiplierless core with
the LOD instruction has \textbf{35.1\% less area}, a \textbf{10\% shorter critical path} (so it can
be clocked faster) and \textbf{17.1\% less power} than the same core with a single-cycle
multiplier. For a tiny battery-powered device, that is a lot.

\begin{kid}
A hardware multiplier is like a big, expensive kitchen machine. Many small kitchens don't
have room for it, so the cook does the job by hand. Hand work is slow. This project is about
finding a quick way of doing it by hand that is \emph{almost} as good as the machine.
\end{kid}

\subsection{The leading one, and the LOD instruction}\label{sec:lod}

The \textbf{leading one} of a number is its highest bit that is 1. Its \emph{position} (counting
from 0 at the right) tells you which power of two the number lies between:

\begin{center}
\begin{tabular}{rlcl}
\toprule
number & bits & leading-one position & so the number is between\\
\midrule
1 & \texttt{1} & 0 & $2^0=1$ and 2\\
13 & \texttt{1101} & 3 & $2^3=8$ and 16\\
11 & \texttt{1011} & 3 & 8 and 16\\
100 & \texttt{1100100} & 6 & 64 and 128\\
200 & \texttt{11001000} & 7 & 128 and 256\\
\bottomrule
\end{tabular}
\end{center}

\textbf{LOD} (``leading-one detector'') is the new instruction from the paper:
\[
  \texttt{lod rd, rs1}\qquad\Longrightarrow\qquad \texttt{rd} \gets \text{position of the leading 1 of } \texttt{rs1}.
\]
For example, if register \code{a1} holds 13, then \code{lod a0, a1} puts 3 into \code{a0}.
Then \code{1 << 3} $=8$ is the largest power of two that is not bigger than 13.

\paragraph{Why not do it in software?} Without the instruction, finding the leading one needs a
loop or a sequence of tests and shifts. On our core, the compiler's software version
(\code{\_\_clzsi2} in libgcc) makes RoBA-Dn1 cost \textbf{91 cycles} per multiply, which is slower
than the exact 84 cycles. With the LOD instruction it costs \textbf{28 cycles}
(Section~\ref{sec:multbench}). Without the new instruction, the whole idea does not pay off.

\paragraph{Why is it cheap in hardware?} In hardware, LOD is a \emph{priority encoder}: a small
circuit that looks at all 32 bits at once and outputs the position of the highest 1. It needs no
memory and no clock cycles of its own. It fits inside the existing arithmetic unit and answers in one
cycle, like an addition. The paper measured that it adds only \textbf{4.5\% power} to the
multiplierless core. Section~\ref{sec:lodpatch} shows the exact lines we added to Ibex.

\begin{keyidea}
LOD answers ``roughly how big is this number?'' in one cycle. Knowing the nearest power of two
of each number turns a multiplication into a few shifts and additions.
\end{keyidea}
\newpage
